\chapter{Theoretical Background}

\section{Reference Frames in Flight Mechanics}

Aircraft motion is described using several coordinate frames \cite{stevens2015aircraft,cook2012flight}. The navigation frame used in this project is a local North-East-Down frame. Its axes point north, east and down, respectively. Altitude is therefore related to the down coordinate by $h=-D$.

The body frame is fixed to the aircraft. Its $x_b$ axis points forward through the nose, its $y_b$ axis points through the right wing and its $z_b$ axis points downward. This convention is right-handed and is compatible with many aerospace simulation formulations. The wind frame is aligned with the air-relative velocity vector and is convenient for defining lift, drag, angle of attack and sideslip.

\section{Six-Degree-of-Freedom Aircraft Motion}

A six-degree-of-freedom aircraft model describes translational and rotational motion of a rigid body. The translational dynamics are obtained from Newton's second law expressed in a rotating body frame. The rotational dynamics are obtained from Euler's rigid-body equations. The state typically contains position, velocity, attitude and angular rates. In this project the state is extended by engine thrust and actuator states so that propulsion and control-surface dynamics are not assumed to be instantaneous.

\section{Aerodynamic Coefficient Modelling}

Aerodynamic forces and moments are commonly represented using nondimensional coefficients. Dynamic pressure is defined as
\begin{equation}
    q_{\infty} = \frac{1}{2}\rho V^2,
\end{equation}
where $\rho$ is air density and $V$ is true airspeed. Lift, drag and moments are then scaled using reference area, span and mean aerodynamic chord. For example,
\begin{equation}
    L = q_{\infty} S C_L, \qquad
    D = q_{\infty} S C_D, \qquad
    M = q_{\infty} S \bar{c} C_m .
\end{equation}

For an engineering first cut, the coefficients may be approximated by linear derivatives and simple nonlinear drag terms. Such a model is not sufficient for high-fidelity prediction near stall or in complex configurations, but it is adequate for early integration and control-design experiments when its validity range is clearly stated.

\section{Propulsion and Actuator Dynamics}

The propulsion model maps throttle commands to thrust. A realistic engine does not change thrust instantaneously, therefore the simulation model should include at least first-order spool dynamics. Similarly, control surfaces should include position limits, rate limits and possibly first-order actuator dynamics. These elements are important because NMPC relies on accurate constraint representation.

\section{Nonlinear Model Predictive Control}

Nonlinear Model Predictive Control solves an optimisation problem at each control update. The controller predicts future system behaviour over a finite horizon and selects an input sequence that minimises a cost function while satisfying constraints \cite{rawlings2017mpc,grune2017nmpc}. Only the first input is applied to the plant, and the optimisation is repeated at the next sampling instant.

NMPC is attractive for aircraft autopilot design because it can explicitly include actuator limits, thrust limits, attitude limits, airspeed bounds, load-factor constraints and path-tracking objectives. Its main disadvantage is computational complexity, which motivates the use of reduced-order prediction models and carefully selected sampling times.
